% ============================================================
% Chapter 6: Experimental Design
%   Section: Evaluation Harness
%     Subsection: Record Format and Run Layout
% Included from chapters/06-experimental-design/04-evaluation-harness/evaluation-harness.tex
% Source: docs/decisions/evaluation.md, "Record format", "Run layout,
% cache and resume"; record.py, store.py.
% ============================================================

\subsection{Record Format and Run Layout}
\label{subsec:harness-records}

\paragraph{One record per claim.} Each claim's result is one line of a
JSON Lines file, with every field present even on an error, so a report
never has to ask whether a key exists. A record holds the gold row's
fields, so a report needs nothing else; the verdict and confidence
after recalibration and the model's raw answer before it; the stage the
claim reached and whether its search or model failed; the queries sent,
with their kind; \emph{every} candidate source seen, once, with the
stage that cut it and why; the ranked evidence in the order the model
saw it, with each factor of its score, its stance, whether it was cited
and the quote; the full trace of progress events with timestamps; the
latency of each stage; the model calls with their tokens; and the
provenance --- model, provider, thresholds, harness version and code
commit. No scraped page body is stored.

\paragraph{Runs are keyed, so the same command resumes.} A run lives in
a directory named by a key derived from the dataset's hash, the model,
the effective thresholds, the scoring weights, the corpus mode and the
harness version: changing any of them starts a new run, while the same
command continues the old one. The commit is recorded on every record
but not keyed, which would discard a run on every commit. Editing a
label changes the dataset's hash, so two gold sets are never mixed in
one run.

\paragraph{Resuming.} Claims already scored are not run again; errors
are retried; a torn last line from a crash is cut and its claim re-run;
and claims whose search or model was unreachable can be re-run on
request, since those are outages, not results. Every record is flushed
and synced to disk as it is written. During the X-Fact pilot this rule
absorbed an interruption of the container runtime and two native
crashes, re-running exactly the claims each had lost
(Section~\ref{sec:results-pilot}). Each session of a run records its
own settings --- timeouts and concurrency --- a fix the pilot prompted.

\paragraph{What is committed.} Raw runs are not committed: they hold
third-party titles and quotes. The report of a run --- its metrics in
JSON, a readable report and a copy of the run's manifest --- is
committed, and those are the numbers this document quotes.
